\chapter{Toxicology, Poisoning, and Environmental Medicine}
\label{ch:toxicology}
Toxicology is the study of adverse effects of chemical, physical, or biological agents on living organisms. Poisoning---accidental or intentional---accounts for millions of emergency department visits and thousands of deaths annually worldwide. Environmental medicine expands this scope to include occupational and environmental exposures. This chapter covers the general approach to the poisoned patient, pharmaceutical overdoses, drugs of abuse, heavy metal toxicity, environmental and occupational exposures, and envenomations.

%-------------------------------------------------------------
\section{Drugs of Abuse}
%-------------------------------------------------------------

%-------------------------------------------------------------

%-------------------------------------------------------------

%-------------------------------------------------------------

Drugs of abuse encompass psychoactive substances capable of inducing chemical dependency, neurotoxicity, and acute toxidrome syndromes. Repeated administration alters central dopaminergic reward pathways and neurotransmitter receptor sensitivity, driving compulsive drug seeking. Management of acute toxicity involves toxidrome recognition, specific antidote administration, intensive supportive care, and long-term addiction treatment.

%-------------------------------------------------------------

%-------------------------------------------------------------

%-------------------------------------------------------------

%-------------------------------------------------------------

\label{sec:Drugs of Abuse}
%-------------------------------------------------------------%-------------------------------------------------------------

\subsection{Alcohol Intoxication and Withdrawal}
\label{sec:Alcohol Intoxication and Withdrawal}
Alcohol (ethanol) is the most widely used psychoactive substance. Intoxication produces dose-dependent CNS depression: disinhibition, ataxia, slurred speech, nystagmus, impaired judgment, and at higher levels, stupor, coma, and respiratory depression. Hypoglycemia (from inhibition of gluconeogenesis) is a critical consideration. Management of intoxication is supportive: monitoring, IV fluids, dextrose (if hypoglycemic), thiamine (100 mg IV) to prevent Wernicke encephalopathy, and consideration of co-ingestants. Alcohol withdrawal syndrome occurs after reduction or cessation of chronic heavy alcohol use, typically starting 6--12 hours after last drink. Clinical features range from mild withdrawal (tremor, anxiety, tachycardia, hypertension) to delirium tremens (profound confusion, agitation, hallucinations, autonomic hyperactivity), with withdrawal seizures (generalized tonic-clonic) occurring 12--48 hours after cessation. Treatment: benzodiazepines are first-line, administered using a symptom-triggered protocol (CIWA-Ar scale); severe DTs may require ICU admission. Prognosis is favorable with appropriate management.

\subsection{Cannabis Toxicity and Cannabinoid Hyperemesis Syndrome}
\label{sec:Cannabis Toxicity}
Cannabis is the most commonly used illicit substance worldwide. Acute intoxication produces euphoria, relaxation, impaired short-term memory, altered judgment, conjunctival injection, tachycardia, and increased appetite. Severe effects include acute psychosis, panic attacks, and in high doses (especially edibles), severe sedation, ataxia, and hypotension. Cannabis hyperemesis syndrome (CHS) is a complication of chronic heavy cannabis use, characterized by cyclic episodes of severe nausea, vomiting, and abdominal pain, relieved by compulsive hot showers or baths. Diagnosis is clinical. Treatment: cessation of cannabis is curative; for acute episodes, hot showers, topical capsaicin cream, benzodiazepines, haloperidol, metoclopramide, and ondansetron. Prognosis is favorable with cannabis abstinence.

\subsection{Hallucinogen Toxicity}
\label{sec:Hallucinogen Toxicity}
LSD (lysergic acid diethylamide) and psilocybin are serotonin 5-HT\textsubscript{2A} receptor agonists producing altered perception, synesthesias, depersonalization, emotional lability, and hallucinations (typically visual). Physiological effects include dilated pupils, tachycardia, hypertension, hyperthermia, and diaphoresis. Acute adverse events include bad trips (anxiety, panic, dysphoria), psychosis, flashbacks, and serotonin syndrome. Phencyclidine (PCP) is a dissociative anesthetic acting as an NMDA receptor antagonist. Toxicity produces nystagmus (horizontal, vertical, or rotatory---diagnostic clue), hypertension, tachycardia, hyperthermia, psychosis, violent behavior, rhabdomyolysis, and seizures. Treatment: benzodiazepines for agitation, antipsychotics for psychosis (caution: haloperidol may lower seizure threshold), and supportive care. Large-volume fluid resuscitation for rhabdomyolysis.

\subsection{MDMA (Ecstasy) Toxicity}
\label{sec:MDMA Toxicity}
MDMA (3,4-methylenedioxymethamphetamine) is a recreational empathogen with stimulant and hallucinogenic properties. It promotes serotonin release and inhibits reuptake. Clinical features include euphoria, increased energy, empathy, and sensory enhancement, but also hyperthermia (often severe), hyponatremia (from SIADH plus excessive water intake), seizures, rhabdomyolysis, DIC, acute liver failure, and serotonin syndrome. Hyponatremia can be severe (Na <120 mEq/L) and cause cerebral edema. Treatment: supportive care---cooling for hyperthermia, hypertonic saline for severe hyponatremia with neurologic symptoms, benzodiazepines for agitation and seizures, cyproheptadine for serotonin syndrome. Prognosis is favorable with appropriate treatment.

\subsection{Novel Synthetic Drugs}
\label{sec:Novel Synthetic Drugs}
Synthetic cannabinoids (Spice, K2) are full agonists at CB1 and CB2 receptors with higher potency and longer half-life than THC. Toxicity: severe agitation, psychosis, seizures, hyperthermia, acute kidney injury, myocardial infarction, and stroke. Synthetic cathinones (bath salts; MDPV, methylone, mephedrone, alpha-PVP) are potent dopamine, norepinephrine, and serotonin reuptake inhibitors and releasers. Toxicity: severe agitation, psychosis, violent behavior, hyperthermia, hyponatremia, serotonin syndrome, rhabdomyolysis, DIC, and cardiovascular collapse. Fentanyl analogs (carfentanil, acetylfentanyl, furanylfentanyl) are ultra-potent opioids, up to 10,000 times more potent than morphine. Toxicity: rapid-onset respiratory depression and death. Treatment: aggressive sedation with benzodiazepines, active cooling, and supportive care; large doses of naloxone may be required for fentanyl analogs. Prognosis is variable.

\subsection{Stimulant Overdose (Cocaine, Methamphetamine)}
\label{sec:Stimulant Overdose}
Cocaine is a dopamine, norepinephrine, and serotonin reuptake inhibitor; methamphetamine promotes release of these monoamines. Both produce a sympathetic storm: hypertension, tachycardia, mydriasis, hyperthermia, psychomotor agitation, psychosis, coronary vasospasm, myocardial infarction, aortic dissection, stroke, rhabdomyolysis, acute kidney injury, and seizures. Cocaine also has sodium channel blockade causing QRS widening. Hyperthermia is the strongest predictor of mortality. Treatment: benzodiazepines (high-dose IV) for agitation and sympathetic hyperactivity, active cooling for hyperthermia, calcium channel blockers or phentolamine for refractory hypertension, and nitroglycerin for cocaine-induced coronary vasospasm. Beta-blockers should be avoided in acute cocaine toxicity due to unopposed alpha-adrenergic stimulation. Prognosis is variable.

%-------------------------------------------------------------
\section{Envenomation and Bites}
%-------------------------------------------------------------

%-------------------------------------------------------------

%-------------------------------------------------------------

%-------------------------------------------------------------

%-------------------------------------------------------------

%-------------------------------------------------------------

%-------------------------------------------------------------

%-------------------------------------------------------------

\label{sec:Envenomation and Bites}
%-------------------------------------------------------------%-------------------------------------------------------------

\subsection{Hymenoptera Stings}
\label{sec:Hymenoptera Stings}
Hymenoptera (bees, wasps, hornets, fire ants) cause local pain, erythema, and swelling from venom components. The major concern is immediate hypersensitivity (IgE-mediated anaphylaxis) in sensitized individuals. Anaphylaxis presents with urticaria, angioedema, bronchospasm, laryngeal edema, hypotension, and cardiovascular collapse occurring within minutes to 1 hour after sting. Diagnosis is clinical. Treatment for anaphylaxis: intramuscular epinephrine (0.3--0.5 mg, 1:1,000, anterolateral thigh) is first-line; adjunctive includes antihistamines, corticosteroids, and IV fluids. Venom immunotherapy (VIT) for patients with a history of anaphylaxis provides 80--98\% protection from future reactions. Prognosis is favorable with prompt treatment.

\subsection{Mammalian Bites}
\label{sec:Mammalian Bites}
Mammalian bites (dog, cat, human) are common and carry risk of infection. Dog bites cause significant crush injury and laceration; pathogens include Pasteurella canis, Staphylococcus aureus, Streptococcus species, and anaerobes. Cat bites create deep puncture wounds carrying bacteria into closed spaces; pathogens include Pasteurella multocida, Staphylococcus, Streptococcus, and Bartonella henselae. Human bites carry high risk of infection from oral flora including Eikenella corrodens, Streptococcus, Staphylococcus, and anaerobes. Management: wound irrigation, debridement of devitalized tissue, assessment of deep structure involvement; primary closure is generally avoided for high-risk wounds except facial wounds. Antibiotic prophylaxis for high-risk wounds. Rabies prophylaxis indicated for bites from unvaccinated or unknown-source wild animals. Tetanus prophylaxis as indicated. Prognosis is favorable with appropriate wound management.

\subsection{Marine Envenomation}
\label{sec:Marine Envenomation}
Jellyfish envenomation: nematocysts discharge venom containing toxins causing local pain, erythema, and in severe cases (box jellyfish---Chironex fleckeri), cardiovascular collapse and death within minutes. Treatment: rinse with vinegar (5\% acetic acid) to inactivate undischarged nematocysts; remove tentacles with forceps; hot water immersion (45\textdegree C, 20 minutes) inactivates venom. For box jellyfish, antivenom is used for severe cases. Stonefish (Synanceia) envenomation produces excruciating pain, local swelling, hypotension, arrhythmias, and rarely death. Treatment: hot water immersion (45\textdegree C for 30--90 minutes denatures the heat-labile toxin); antivenom is available. Cone snail (Conus) venom contains conotoxins causing neuromuscular blockade, bulbar palsy, and respiratory failure. Treatment is supportive with airway support and ventilation as needed. Prognosis is favorable with prompt treatment.

\subsection{Scorpion Stings}
\label{sec:Scorpion Stings}
Centruroides sculpturatus (bark scorpion) is the only medically significant scorpion in North America. Venom is a sodium channel toxin causing prolonged neuronal depolarization and excessive autonomic and somatic motor activity. Clinical features: immediate local pain, paresthesias, involuntary muscle spasms, cranial nerve dysfunction (nystagmus, opsoclonus, hypersalivation, tongue fasciculations, dysphagia), respiratory distress, and seizures. Children are at highest risk for severe toxicity. Diagnosis is clinical. Management: supportive care, benzodiazepines for agitation and muscle spasms. Anascorp (Centruroides immune F(ab')2 antivenom) is highly effective, with resolution of symptoms within 1--4 hours. Prognosis is favorable with antivenom administration.

\subsection{Snake Envenomation}
\label{sec:Snake Envenomation}
Clinically significant snake envenomation is classified by family: Viperidae (vipers, pit vipers---rattlesnakes, copperheads, cottonmouths) and Elapidae (cobras, kraits, mambas, coral snakes, sea snakes). Viperidae venom contains hemotoxins causing local tissue destruction, coagulopathy, and thrombocytopenia. Elapidae venom contains neurotoxins causing neuromuscular blockade with descending paralysis, ptosis, bulbar palsy, respiratory failure, and sometimes minimal local symptoms. Sea snakes cause myotoxicity with generalized myalgia, myoglobinuria, and acute kidney injury. Clinical evaluation uses crotalid envenomation severity grading. Diagnosis is clinical with lab evaluation (CBC, PT/PTT, INR, fibrinogen, D-dimer, creatinine, CK). Management: supportive care, wound management, tetanus prophylaxis, and antivenom (Crotalidae polyvalent immune Fab for North American pit vipers). Prognosis is favorable with timely antivenom administration.

\subsection{Spider Bites}
\label{sec:Spider Bites}
Clinically significant spider bites in North America include black widow (Latrodectus) and brown recluse (Loxosceles). Black widow venom contains alpha-latrotoxin, causing massive release of acetylcholine and norepinephrine. Clinical features: local pinprick sensation, then severe muscle pain and cramping, rigidity, diaphoresis, nausea, vomiting, hypertension, tachycardia, and rarely respiratory compromise. Treatment: supportive care, benzodiazepines for muscle spasms, and latrodectus antivenom for severe cases. Brown recluse venom contains sphingomyelinase D, causing dermonecrotic injury, hemolysis, and coagulopathy. Clinical features: initially painless bite, then over 6--12 hours erythema, blister, and necrosis forming an eschar over 1--2 weeks. Systemic loxoscelism includes fever, chills, nausea, arthralgias, hemolytic anemia, thrombocytopenia, DIC, acute kidney injury. Treatment: wound care, supportive care for hemolysis and renal failure. Prognosis is generally favorable.

%-------------------------------------------------------------
\section{Environmental and Occupational Medicine}
%-------------------------------------------------------------

%-------------------------------------------------------------

%-------------------------------------------------------------

%-------------------------------------------------------------

%-------------------------------------------------------------

%-------------------------------------------------------------

%-------------------------------------------------------------

%-------------------------------------------------------------

\label{sec:Environmental and Occupational Medicine}
%-------------------------------------------------------------%-------------------------------------------------------------

\subsection{Altitude Sickness}
\label{sec:Altitude Sickness}
Altitude illness occurs at elevations above 2,500 m (8,000 ft) due to decreased barometric pressure and hypoxia. The condition results from compensatory hyperventilation, increased sympathetic tone, and pulmonary and cerebral vasoconstriction and vasodilation; inadequate acclimatization leads to three syndromes. Acute mountain sickness (AMS) is the most common form: headache, nausea, dizziness, anorexia, fatigue, and sleep disturbance. Diagnosis uses the Lake Louise Score. Treatment: acetazolamide for prophylaxis and mild treatment; descent for persistent symptoms. High-altitude cerebral edema (HACE) is a life-threatening progression with ataxia, altered mental status, hallucinations, seizures, and coma. Treatment: immediate descent, supplemental oxygen, dexamethasone. High-altitude pulmonary edema (HAPE) is non-cardiogenic pulmonary edema with cough, dyspnea at rest, cyanosis, and crackles. Treatment: immediate descent, supplemental oxygen, nifedipine or phosphodiesterase-5 inhibitors. Prognosis is favorable with prompt descent.

\subsection{Carbon Monoxide Poisoning}
\label{sec:Carbon Monoxide Poisoning}
Carbon monoxide (CO) is an odorless, colorless gas produced by incomplete combustion of carbon-based fuels. The condition results from CO binding to hemoglobin with 240 times greater affinity than oxygen, forming carboxyhemoglobin (COHb), which shifts the oxygen dissociation curve to the left, impairing oxygen delivery and utilization; CO also binds to myoglobin and cytochrome c oxidase. Clinical features: headache (most common), nausea, dizziness, confusion, dyspnea, blurred vision, loss of consciousness, and seizures; severe poisoning causes coma, respiratory failure, cardiovascular collapse, and death. Neurologic sequelae include acute encephalopathy and delayed neuropsychiatric syndrome (DNS). Diagnosis is by COHb level. Treatment: 100\% normobaric oxygen; hyperbaric oxygen therapy (HBOT) is indicated for severe poisoning and may prevent DNS. Prognosis is favorable with prompt treatment.

\subsection{Cyanide Toxicity}
\label{sec:Cyanide Toxicity}
Cyanide is a rapidly acting poison from industrial sources, smoke inhalation, and certain plants. The condition results from cyanide binding to cytochrome c oxidase (complex IV of the mitochondrial electron transport chain), inhibiting cellular respiration and causing histotoxic hypoxia. Clinical features: rapid onset (seconds to minutes) of altered mental status, seizures, coma, trismus, hyperventilation, hypotension, bradycardia, ventricular arrhythmias, and cardiovascular collapse; classic sign is bright red venous blood. Diagnosis is clinical, especially in smoke inhalation victims; elevated lactate suggests cyanide toxicity. Management: hydroxocobalamin (5 g IV) is first-line, binding cyanide to form nontoxic cyanocobalamin; sodium thiosulfate provides substrate for rhodanese to convert cyanide to thiocyanate. Prognosis is favorable with rapid treatment.

\subsection{Drowning (Submersion Injury)}
\label{sec:Drowning}
Drowning is respiratory impairment from submersion or immersion in a liquid medium. It is a leading cause of unintentional injury death worldwide, particularly in children under 5 years. The condition results from initial breath-holding and panic, laryngospasm (10--15\% dry drowning---no aspiration), then aspiration of water, surfactant washout, pulmonary edema, and hypoxemia. Clinical features: respiratory distress, coughing, tachypnea, cyanosis, foam from mouth/nose, altered mental status, seizures, cardiac arrest; the primary determinant of outcome is duration of anoxia. Diagnosis is by history of submersion, clinical presentation, chest X-ray, and arterial blood gas. Management: immediate rescue, CPR if pulseless, airway management, positive pressure ventilation, supplemental oxygen, and rewarming for hypothermic victims. Prognosis is better in children and those with submersion time <10 minutes and rapid resuscitation.

\subsection{Electrical Injury and Lightning Strike}
\label{sec:Electrical Injury}
Electrical injuries from high-voltage (>1,000 V) and low-voltage (<1,000 V, typically household 110--240 V) sources cause damage via thermal burns, direct cellular injury, and electroporation of cell membranes. Current follows the path of least resistance: nerves and blood vessels are most affected, followed by muscle, then skin and bone. Clinical features: cutaneous burns (entry/exit wounds---may be deceptively small), muscle necrosis, rhabdomyolysis, compartment syndrome, cardiac arrhythmias, neurological injury, and fractures. Lightning strike produces a massive direct current (DC) shock of up to 30 million volts and 30,000 amps. Clinical features: cardiac arrest, respiratory arrest, confusion, amnesia, kerauonographic markings, tympanic membrane rupture, and cataracts. Diagnosis is by history and physical examination. Management: advanced cardiac life support, fluid resuscitation, surgical fasciotomy for compartment syndrome, and wound care. Prognosis is variable; lightning strike patients in cardiac arrest may be salvageable with prolonged resuscitation.

\subsection{Haff Disease}
\label{sec:Haff Disease}
Haff disease is a rare syndrome of rhabdomyolysis occurring within 24 hours of consuming certain freshwater or marine fish or crustaceans. The etiology is an unknown toxin (possibly a heat-stable palytoxin-like compound from algae/dinoflagellates that accumulates in the food chain). Clinical features: sudden onset of severe myalgia, muscle tenderness, weakness, dark urine (myoglobinuria), and elevated serum creatine kinase (CK, often >10,000 IU/L) within 6--24 hours of fish consumption, without fever, neurologic deficits, or rash. Diagnosis is clinical with CK elevation, urinalysis (heme-positive without RBCs), and history of recent fish consumption. Management: aggressive IV fluids (target urine output >200 mL/hour), monitoring for acute kidney injury and compartment syndrome, and correction of electrolyte abnormalities. Prognosis is excellent; most patients recover within days without long-term sequelae if treated promptly.

\subsection{Harmful Algal Blooms}
\label{sec:Harmful Algal Blooms}
Harmful algal blooms (HABs) are caused by toxin-producing cyanobacteria (blue-green algae) and marine algae that accumulate in freshwater lakes, ponds, and coastal waters during warm weather. Cyanobacteria produce multiple toxin classes: hepatotoxins (microcystins, nodularins), neurotoxins (anatoxin-a, saxitoxin), and dermatoxins (lyngbyatoxin). Exposure occurs through recreational water contact, drinking contaminated water, or consumption of toxin-accumulating shellfish. Clinical features: contact---dermatitis, conjunctivitis, allergic rhinitis; ingestion---acute gastroenteritis, hepatotoxicity, neurotoxicity; inhalation---asthma-like symptoms. Diagnosis is clinical with history of exposure to water with visible algal blooms. Management: supportive---decontamination, IV fluids for GI losses, respiratory support for neurotoxic effects. No specific antidotes exist. Prevention includes avoiding water contact where blooms are visible. Climate change is increasing HAB frequency, duration, and geographic range.

\subsection{Heat Stroke and Hyperthermia}
\label{sec:Heat Stroke}
Heat stroke is a life-threatening condition characterized by core body temperature >40\textdegree C (104\textdegree F) with CNS dysfunction. Two forms: classic (nonexertional)---affects elderly, infants, and those in heat waves; and exertional---affects athletes, military personnel, and laborers. The condition results from failure of thermoregulatory mechanisms with overwhelming heat production or impaired heat dissipation, causing direct cellular injury, endotoxemia, systemic inflammatory response, and multiorgan failure. Clinical features: hyperthermia, anhydrosis in classic heat stroke (skin hot and dry), CNS dysfunction (confusion, ataxia, delirium, seizures, coma), tachycardia, hypotension, and circulatory collapse. Diagnosis: core temperature >40\textdegree C with CNS dysfunction after ruling out other causes. Management: rapid cooling---evaporative cooling, ice packs, cold IV fluids, and ice water immersion; aggressive IV fluids; benzodiazepines for uncontrolled shivering or seizures. Prognosis depends on promptness of treatment.

\subsection{Hydrogen Sulfide Toxicity}
\label{sec:Hydrogen Sulfide Toxicity}
Hydrogen sulfide (H\textsubscript{2}S) is a colorless gas with a characteristic rotten egg odor (rapid olfactory fatigue occurs, rendering odor unreliable), encountered in sewers, manure pits, oil and gas drilling, and pulp mills. The condition results from H\textsubscript{2}S inhibiting cytochrome c oxidase (identical mechanism to cyanide), causing histotoxic hypoxia. Clinical features: acute exposure causes conjunctivitis, respiratory tract irritation, pulmonary edema, loss of consciousness (knockdown---immediate collapse), seizures, and death; chronic low-level exposure causes headache, fatigue, olfactory dysfunction, and cognitive impairment. Diagnosis is clinical. Management: immediate removal from exposure, nitrite therapy to induce methemoglobinemia, hyperbaric oxygen, and supportive care. Prognosis depends on severity of exposure.

\subsection{Hypothermia}
\label{sec:Hypothermia}
Hypothermia is defined as core temperature below 35\textdegree C (95\textdegree F), classified as mild (32--35\textdegree C), moderate (28--32\textdegree C), and severe (<28\textdegree C). Causes include environmental exposure, cold-water immersion, alcohol intoxication, hypothyroidism, hypoglycemia, and trauma. The condition results from initial sympathetic activation (shivering, vasoconstriction, tachycardia) transitioning to progressive depression of all organ systems as temperature drops. Clinical features: mild hypothermia: intact shivering, confusion, ataxia, dysarthria; moderate hypothermia: cessation of shivering, progressive bradycardia, hyporeflexia, paradoxical undressing, obtundation; severe hypothermia: coma, apnea, ventricular arrhythmias, asystole. The classic ECG finding is the Osborn wave (J wave). Diagnosis is by core temperature measurement. Management: passive external rewarming for mild hypothermia; active external rewarming for moderate hypothermia; active internal rewarming (warmed IV fluids, ECMO) for severe hypothermia. Prognosis is favorable with appropriate rewarming.

\subsection{Methyl Parathion and Other Organophosphate Insecticides}
\label{sec:Methyl Parathion}
Methyl parathion is a highly toxic organophosphate insecticide (WHO class Ia---extremely hazardous) used in agriculture. Like other organophosphates, it irreversibly inhibits acetylcholinesterase, causing cholinergic crisis (SLUDGE syndrome plus miosis, bronchorrhea, bronchospasm, bradycardia, muscle fasciculations, weakness, seizures, coma). It is particularly dangerous because it is rapidly absorbed through skin. Diagnosis is clinical with low RBC/plasma cholinesterase activity (<50\% of normal). Treatment: atropine (to block muscarinic symptoms), pralidoxime (2-PAM, to reactivate AChE), diazepam for seizures, and decontamination. Prognosis: recovery is usually complete if treated early; delayed neurotoxicity (organophosphate-induced delayed neuropathy/OPIDN) can occur weeks later.

\subsection{Organophosphate and Carbamate Insecticide Poisoning}
\label{sec:Organophosphate Poisoning}
Organophosphate (OP) and carbamate insecticides inhibit acetylcholinesterase (AChE) at cholinergic synapses, causing accumulation of acetylcholine and excessive muscarinic, nicotinic, and CNS stimulation. OPs bind irreversibly (aging occurs within 24--48 hours); carbamates bind reversibly. Clinical features: cholinergic crisis---SLUDGE syndrome (Salivation, Lacrimation, Urination, Defecation, GI upset, Emesis) plus miosis, bronchospasm, bronchorrhea, bradycardia, excessive sweating, muscle fasciculations, weakness, paralysis, seizures, and CNS depression. Diagnosis is clinical with RBC AChE activity for confirmation. Management: immediate decontamination, atropine (1--4 mg IV, doubled every 5--10 minutes) to block muscarinic effects, pralidoxime (2-PAM, 1--2 g IV) to reactivate AChE if given before aging, and benzodiazepines for seizures. Prognosis is favorable with early treatment.

\subsection{Radiation Sickness (Acute Radiation Syndrome)}
\label{sec:Radiation Sickness}
Acute radiation syndrome (ARS) results from high-dose whole-body or significant partial-body ionizing radiation exposure (>1 Gy). The condition results from radiation causing DNA damage (double-strand breaks), cell death in rapidly dividing tissues (bone marrow, GI tract, skin), and free-radical generation. Clinical phases: prodromal phase (0--48 hours): nausea, vomiting, diarrhea, anorexia; latent phase: transient symptom improvement; manifest illness phase depends on dose: hematopoietic syndrome (1--6 Gy): pancytopenia, infection, hemorrhage, anemia; gastrointestinal syndrome (6--10 Gy): severe diarrhea, fluid loss, mucosal necrosis, sepsis; cerebrovascular/cardiovascular syndrome (>10--20 Gy): hypotension, cerebral edema, convulsions, coma, death. Diagnosis is based on history of radiation exposure, absolute lymphocyte count decline, and cytogenetic analysis. Treatment: supportive care (reverse isolation, antimicrobial prophylaxis, blood product support), cytokines (G-CSF, GM-CSF) for neutropenia, and treatment of internal contamination. Prognosis is poor for doses >6 Gy.

%-------------------------------------------------------------
\section{General Approach to the Poisoned Patient}
%-------------------------------------------------------------

%-------------------------------------------------------------

%-------------------------------------------------------------

%-------------------------------------------------------------

%-------------------------------------------------------------

%-------------------------------------------------------------

%-------------------------------------------------------------

%-------------------------------------------------------------

\label{sec:General Approach to the Poisoned Patient}
%-------------------------------------------------------------%-------------------------------------------------------------

\subsection{Antidotes Overview}
\label{sec:Antidotes Overview}
Key antidotes include:
\begin{itemize}[noitemsep]
  \item \textbf{Naloxone:} Opioid receptor antagonist; 0.04--2 mg IV/IM/IN for respiratory depression from opioids; may repeat or infuse.
  \item \textbf{N-Acetylcysteine (NAC):} Replenishes hepatic glutathione for acetaminophen toxicity; IV (20-hour protocol) or oral (72-hour protocol) within 8--10 hours.
  \item \textbf{Flumazenil:} GABA\textsubscript{A} receptor antagonist for benzodiazepine overdose; used with extreme caution due to risk of seizures in mixed overdoses or chronic benzodiazepine users.
  \item \textbf{Fomepizole:} Alcohol dehydrogenase inhibitor for methanol and ethylene glycol poisoning; preferred over ethanol.
  \item \textbf{Physostigmine:} Acetylcholinesterase inhibitor for anticholinergic toxidrome (pure anticholinergic delirium, supraventricular tachycardia); contraindicated in TCA overdose.
  \item \textbf{Digoxin Immune Fab (Digibind):} Antibody fragments that bind digoxin for life-threatening digoxin toxicity.
  \item \textbf{Glucagon:} Positive chronotrope/inotrope via cAMP activation for beta-blocker and calcium channel blocker overdose.
  \item \textbf{Hydroxocobalamin:} Vitamin B\textsubscript{12a} that binds cyanide; first-line antidote for cyanide poisoning.
  \item \textbf{Atropine:} Antimuscarinic for organophosphate poisoning and symptomatic bradycardia.
  \item \textbf{Pralidoxime (2-PAM):} Reactivates acetylcholinesterase in organophosphate poisoning.
  \item \textbf{Octreotide:} Somatostatin analog used for hypoglycemia from sulfonylurea overdose.
  \item \textbf{Sodium bicarbonate:} Alkalinizes serum to treat TCA-induced QRS widening and salicylate toxicity.
\end{itemize}

\subsection{Initial Stabilization and Decontamination}
\label{sec:Initial Stabilization}
The poisoned patient is managed using an ABCDE (Airway, Breathing, Circulation, Disability, Exposure/Decontamination) approach. Airway assessment is critical in patients with decreased level of consciousness (GCS <8 often requires intubation). Breathing support with supplemental oxygen and monitoring of pulse oximetry and end-tidal CO\textsubscript{2} is essential for agents causing respiratory depression. Circulation involves IV access, continuous cardiac monitoring, blood pressure support with IV fluids, and vasopressors if needed. Disability assessment includes mental status (AVPU/GCS), pupil size, and glucose measurement. Decontamination includes removal of contaminated clothing, skin decontamination with soap and water, and gastrointestinal decontamination with activated charcoal (1 g/kg, within 1--2 hours of ingestion for most toxins). Gastric lavage is rarely indicated (only for life-threatening ingestions within 1 hour). Whole-bowel irrigation may be used for sustained-release preparations, iron, lithium, and body packers.

\subsection{Toxicokinetics and Toxicodynamics}
\label{sec:Toxicokinetics and Toxicodynamics}
Toxicokinetics describes the absorption, distribution, metabolism, and elimination of toxins. Absorption occurs via ingestion (most common), inhalation, dermal contact, or injection. Distribution depends on protein binding, lipid solubility, and volume of distribution (V\textsubscript{d}). Metabolism primarily occurs in the liver via Phase I (cytochrome P450 oxidation, reduction, hydrolysis) and Phase II (glucuronidation, sulfation, acetylation, glutathione conjugation) reactions. Elimination is via renal, hepatic (biliary), or pulmonary routes. Toxicodynamics describes the molecular mechanisms by which toxins produce their effects---receptor binding (opioids, benzodiazepines), enzyme inhibition (organophosphates, cyanide), ion channel blockade (TCA, calcium channel blockers), or mitochondrial toxicity (acetaminophen NAPQI, cyanide).

%-------------------------------------------------------------
\section{Heavy Metal Toxicity}
%-------------------------------------------------------------

%-------------------------------------------------------------

%-------------------------------------------------------------

%-------------------------------------------------------------

%-------------------------------------------------------------

%-------------------------------------------------------------

%-------------------------------------------------------------

%-------------------------------------------------------------

\label{sec:Heavy Metal Toxicity}
%-------------------------------------------------------------%-------------------------------------------------------------

\subsection{Arsenic Poisoning}
\label{sec:Arsenic Poisoning}
Arsenic is a naturally occurring metalloid found in groundwater, rice, and some industrial products. The condition results from inorganic arsenic (trivalent arsenite) inhibiting pyruvate dehydrogenase, disrupting the Kreb's cycle and oxidative phosphorylation, and impairing DNA repair. Acute arsenic poisoning presents with severe gastrointestinal symptoms (vomiting, profuse watery diarrhea), hypotension, cardiac dysfunction, acute kidney injury, and death within 24 hours. Subacute/chronic toxicity presents with peripheral neuropathy, Mee's lines (transverse white bands on fingernails), diffuse hyperpigmentation, palmar-plantar hyperkeratosis, and squamous cell carcinoma. Diagnosis is by 24-hour urine arsenic level. Treatment: dimercaprol (BAL) or succimer (DMSA) for acute symptomatic poisoning. Prognosis depends on severity and timing of treatment.

\subsection{Cadmium Toxicity}
\label{sec:Cadmium Toxicity}
Cadmium is a heavy metal used in batteries, pigments, electroplating, and fertilizers, with a biological half-life of 15--30 years. The condition results from cadmium accumulating in the kidneys and bone, impairing renal tubular reabsorption and causing osteomalacia. Acute inhalation exposure causes chemical pneumonitis and pulmonary edema. Chronic exposure causes renal tubular damage, Itai-Itai disease (severe osteomalacia, bone pain, and fractures with renal tubular dysfunction), hypertension, and lung cancer. Diagnosis: 24-hour urine and blood cadmium levels. Treatment: removal from exposure; chelation is generally ineffective and BAL is contraindicated. Prognosis depends on chronicity of exposure.

\subsection{Iron Overdose}
\label{sec:Iron Overdose}
Iron overdose is a leading cause of poisoning death in children under 6 years due to accidental ingestion of prenatal iron supplements. Toxic dose: elemental iron >20 mg/kg; severe toxicity >60 mg/kg. The condition results from free iron catalyzing hydroxyl radical formation via the Fenton reaction, causing lipid peroxidation, mitochondrial dysfunction, and metabolic acidosis. Clinical stages: Stage 1 (0--6 hours): nausea, vomiting, abdominal pain, diarrhea, hematemesis, hematochezia; Stage 2 (6--24 hours): apparent recovery; Stage 3 (12--48 hours): metabolic acidosis, seizures, coma, shock, coagulopathy, hepatic necrosis, pulmonary edema; Stage 4 (2--4 weeks): gastrointestinal obstruction from strictures. Diagnosis is by serum iron level and abdominal X-ray. Treatment: aggressive IV fluids, deferoxamine (IV iron chelator) for metabolic acidosis, shock, or iron level >500 mcg/dL, and whole-bowel irrigation for retained tablets. Prognosis depends on severity.

\subsection{Lead Poisoning (Plumbism)}
\label{sec:Lead Poisoning}
Lead poisoning (plumbism) is a cumulative toxicant affecting multiple organ systems. Sources include lead-based paint (most common in children), contaminated water, soil, industrial exposure, traditional cosmetics, and imported products. The condition results from lead impairing heme synthesis, disrupting calcium homeostasis, and causing oxidative stress; in children, lead crosses the blood-brain barrier and impairs neuronal migration, synaptogenesis, and myelination. No safe blood lead level exists; the CDC reference level is 3.5 mcg/dL. Clinical features in children include asymptomatic (most cases), neurocognitive deficits, developmental delay, and at very high levels: encephalopathy, seizures, coma, and death. In adults: peripheral neuropathy, abdominal pain, arthralgias, fatigue, headache, hypertension, nephropathy, and anemia. Diagnosis is by venous blood lead level. Treatment includes removal from exposure source and chelation therapy: succimer (DMSA, oral) for asymptomatic children with levels $\ge$45 mcg/dL; BAL plus calcium disodium EDTA for severe toxicity. Prognosis is favorable with early intervention.

\subsection{Mercury Toxicity}
\label{sec:Mercury Toxicity}
Mercury exists in three forms: elemental (metallic), inorganic (mercurous/mercuric salts), and organic (methylmercury). Elemental mercury: inhaled vapor is absorbed (80\%); causes acute pneumonitis, pulmonary edema, and chronic neurotoxicity (tremor, erethism---emotional lability, pathological shyness, memory loss). Inorganic mercury: ingestion of mercuric chloride causes corrosive gastroenteritis, acute tubular necrosis, and nephrotic syndrome. Acrodynia (Pink disease) is a hypersensitivity reaction in children. Organic mercury (methylmercury): bioaccumulates in fish; causes Minamata disease (severe neurological syndrome: paresthesias, ataxia, dysarthria, visual field constriction, hearing loss) and is a potent teratogen. Diagnosis: blood and urine mercury levels. Treatment: cessation of exposure; succimer (DMSA) is the preferred chelator for elemental and inorganic mercury. Prognosis depends on form and severity of exposure.

%-------------------------------------------------------------
\section{Pharmaceutical Overdoses}
%-------------------------------------------------------------

%-------------------------------------------------------------

%-------------------------------------------------------------

%-------------------------------------------------------------

%-------------------------------------------------------------

%-------------------------------------------------------------

%-------------------------------------------------------------

%-------------------------------------------------------------

\label{sec:Pharmaceutical Overdoses}
%-------------------------------------------------------------%-------------------------------------------------------------

\subsection{Acetaminophen (Paracetamol) Overdose}
\label{sec:Acetaminophen Overdose}
Acetaminophen overdose is the most common pharmaceutical overdose in the United States and the leading cause of acute liver failure. A toxic dose in adults is typically >7.5--10 g (or >150 mg/kg). The condition results from depletion of hepatic glutathione and accumulation of the toxic metabolite N-acetyl-p-benzoquinone imine (NAPQI), which covalently binds to hepatocyte proteins, causing centrilobular hepatic necrosis. Clinical features progress through four phases: Phase I (0--24 hours): nausea, vomiting, malaise, pallor; Phase II (24--72 hours): right upper quadrant pain, elevated transaminases, PT/INR prolongation; Phase III (72--96 hours): peak hepatotoxicity, hepatic encephalopathy, coagulopathy, jaundice, acute liver failure; Phase IV (4 days--2 weeks): recovery or progression to death. Diagnosis is by serum acetaminophen level at >4 hours post-ingestion plotted on the Rumack-Matthew nomogram. Management involves N-acetylcysteine (NAC), which replenishes glutathione and provides substrate for non-toxic sulfate conjugation, most effective if started within 8--10 hours. Prognosis depends on timely treatment; the King's College Criteria guide liver transplantation decisions.

\subsection{Benzodiazepine Overdose}
\label{sec:Benzodiazepine Overdose}
Benzodiazepine overdose is a pharmaceutical overdose involving positive allosteric modulators of GABA\textsubscript{A} receptors. The condition results from excessive GABAergic inhibition producing CNS depression. Clinical features include drowsiness, slurred speech, ataxia, and coma; respiratory depression is mild in pure benzodiazepine overdose but is synergistic with ethanol, opioids, and other CNS depressants. Diagnosis is clinical; urine drug screens have limited utility. Management is supportive care (airway protection, respiratory support); flumazenil, a GABA\textsubscript{A} receptor antagonist, is reserved for pure benzodiazepine overdose with severe respiratory depression and is contraindicated in chronic benzodiazepine users or mixed ingestions due to risk of seizures. Prognosis is favorable with supportive care.

\subsection{Beta-Blocker Overdose}
\label{sec:Beta-Blocker Overdose}
Beta-blocker overdose (propranolol most lethal, also atenolol, metoprolol, sotalol) is a pharmaceutical overdose causing bradycardia, hypotension, and decreased myocardial contractility. Propranolol also has sodium channel blockade contributing to QRS widening and seizures; sotalol causes QT prolongation and torsades de pointes. The condition results from beta-adrenergic receptor blockade. Clinical features include bradyarrhythmias, low cardiac output, pulmonary edema, hypoglycemia, seizures, and altered mental status. Diagnosis is clinical. Management includes IV fluids, atropine, glucagon (5--10 mg IV bolus, then 1--10 mg/hour infusion) which increases cyclic AMP via non-adrenergic activation of adenylate cyclase, high-dose insulin euglycemia therapy, vasopressors, and calcium salts. Prognosis is variable.

\subsection{Calcium Channel Blocker Overdose}
\label{sec:CCB Overdose}
Calcium channel blocker (CCB) overdose (verapamil, diltiazem, nifedipine, amlodipine, and other dihydropyridines) is a highly lethal pharmaceutical overdose. The condition results from CCBs inhibiting L-type calcium channels in cardiac myocytes, vascular smooth muscle, and pancreatic beta cells, leading to negative chronotropy, dromotropy, and inotropy with vasodilation. Clinical features include hypotension, bradycardia (or reflex tachycardia with dihydropyridines), atrioventricular block, junctional rhythms, asystole, hyperglycemia, metabolic acidosis, and altered mental status. Diagnosis is clinical. Management includes aggressive fluid resuscitation, calcium salts to overcome calcium channel blockade, and high-dose insulin euglycemia therapy (the most effective treatment). Prognosis is guarded; vasopressors and ECMO may be needed.

\subsection{Digoxin Toxicity}
\label{sec:Digoxin Toxicity}
Digoxin toxicity is a pharmaceutical overdose involving a cardiac glycoside that inhibits Na+/K+-ATPase, increasing intracellular calcium and enhancing cardiac contractility. Toxicity occurs with therapeutic excess (most common, especially with renal impairment, hypokalemia, hypomagnesemia, hypercalcemia, or drug interactions) or acute overdose. Clinical features include gastrointestinal (nausea, vomiting, anorexia), neurological (visual disturbances---yellow-green halos, blurred vision, confusion), and cardiac manifestations including any bradyarrhythmia, atrial tachycardia with block, accelerated junctional rhythm, and bidirectional ventricular tachycardia (pathognomonic); hyperkalemia is a marker of severe acute toxicity. Diagnosis is clinical with ECG findings and serum digoxin level. Management includes digoxin immune Fab (Digibind) as the definitive antidote for life-threatening arrhythmias or hyperkalemia. Prognosis is favorable with timely administration of Fab fragments.

\subsection{Lithium Toxicity}
\label{sec:Lithium Toxicity}
Lithium toxicity is a pharmaceutical overdose involving a monovalent cation used for bipolar disorder with a narrow therapeutic index (0.6--1.2 mEq/L). Toxicity occurs from acute ingestion, acute-on-chronic overdose, or chronic accumulation (often due to dehydration, renal impairment, or drug interactions). The condition results from lithium substituting for sodium in excitable cells, altering neurotransmitter release, and impairing renal concentrating ability. Clinical features: acute toxicity presents with gastrointestinal symptoms and mild neurologic symptoms (tremor, ataxia, nystagmus); chronic toxicity presents predominantly with neurotoxicity (confusion, dysarthria, ataxia, hyperreflexia, clonus, parkinsonism), with severe cases progressing to seizures, coma, and permanent cerebellar injury. Diagnosis is by serum lithium level (interpreted with caution---chronic toxicity can occur at therapeutic levels). Management includes IV fluids (normal saline) to correct volume depletion and enhance lithium clearance; hemodialysis is indicated for severe toxicity. Prognosis is favorable with prompt treatment; chronic toxicity may cause permanent neurologic injury.

\subsection{Metformin Overdose}
\label{sec:Metformin Overdose}
Metformin overdose is an uncommon but highly lethal pharmaceutical overdose due to metformin-associated lactic acidosis (MALA). The condition results from metformin inhibiting gluconeogenesis and complex I of the mitochondrial electron transport chain, leading to impaired lactate clearance and increased anaerobic metabolism, with predisposing factors including renal impairment, hepatic dysfunction, hypoperfusion, sepsis, and concurrent use of contrast dye. Clinical features include nausea, vomiting, abdominal pain, diarrhea, profound lactic acidosis (pH <7.1, lactate >20 mmol/L), hypotension, hypothermia, altered mental status, and acute kidney injury. Diagnosis is by elevated serum metformin level and severe lactic acidosis with elevated lactate-to-pyruvate ratio. Management includes supportive care and early hemodialysis or continuous venovenous hemofiltration as definitive treatment. Prognosis is guarded; severe MALA has mortality exceeding 30--50\%.

\subsection{Opioid Overdose}
\label{sec:Opioid Overdose}
Opioid overdose is a leading cause of poisoning death, driven by prescription opioids, heroin, and illicit fentanyl analogs. The condition results from mu-opioid receptor agonism in the brainstem producing respiratory depression, miosis (pinpoint pupils), and CNS depression leading to coma. Clinical features include the classic triad of coma, miosis, and respiratory depression. Diagnosis is clinical. Management involves airway management and naloxone (0.04--2 mg IV, 0.4--2 mg IM, or 2--4 mg IN) as the antidote; naloxone has a short half-life (30--90 minutes) that may require repeat doses or continuous infusion. Patients should be observed for at least 4--6 hours after the last naloxone dose. Prognosis is favorable with timely intervention.

\subsection{Salicylate (Aspirin) Overdose}
\label{sec:Salicylate Overdose}
Salicylate overdose is a pharmaceutical overdose that can be acute or chronic (chronic toxicity more dangerous with higher mortality). The condition results from salicylates uncoupling oxidative phosphorylation, stimulating the respiratory center, inhibiting Kreb's cycle enzymes, and increasing metabolic acid production. Clinical features include tinnitus (earliest symptom), hyperventilation, nausea, vomiting, hyperthermia, diaphoresis, and tachycardia; as toxicity progresses, mixed acid-base disorders emerge (respiratory alkalosis followed by concurrent metabolic acidosis); severe toxicity includes altered mental status, seizures, coma, and pulmonary edema. Diagnosis is by serum salicylate level. Management includes activated charcoal, IV fluids with dextrose, bicarbonate therapy (goal urine pH 7.5--8.0) to enhance renal elimination, and hemodialysis for severe toxicity. Prognosis is favorable with appropriate treatment.

\subsection{SSRI Overdose and Serotonin Syndrome}
\label{sec:SSRI Overdose}
Selective serotonin reuptake inhibitor (SSRI) overdose is generally less toxic than TCA overdose. SSRIs cause mild CNS depression, nausea, vomiting, and tachycardia. Citalopram and escitalopram are associated with QT prolongation (dose-dependent). Serotonin syndrome is a potentially life-threatening condition caused by excessive serotonergic activity, most commonly from SSRI in combination with MAOIs, linezolid, methylene blue, dextromethorphan, tramadol, meperidine, or MDMA. The triad is altered mental status, autonomic hyperactivity, and neuromuscular abnormalities. Clinical features include clonus, hyperreflexia, tremor, myoclonus, hyperthermia, diaphoresis, tachycardia, hypertension, and agitation. Diagnosis is clinical using the Hunter Serotonin Toxicity Criteria. Management includes discontinuation of serotonergic agents, supportive care (cooling, IV fluids, benzodiazepines), and cyproheptadine for moderate to severe cases. Prognosis is favorable with early recognition and treatment.

\subsection{Theophylline Overdose}
\label{sec:Theophylline Overdose}
Theophylline overdose is a pharmaceutical overdose causing severe neurological and cardiovascular toxicity. Theophylline is a phosphodiesterase and adenosine receptor antagonist, leading to increased catecholamine release and beta-adrenergic stimulation. Clinical features: acute toxicity presents with nausea, vomiting, hypokalemia, hypophosphatemia, hyperglycemia, metabolic acidosis, tremor, agitation, and seizures (often refractory); cardiovascular effects include sinus tachycardia, atrial fibrillation, ventricular tachycardia, and hypotension. Diagnosis is by serum theophylline level. Management includes multidose activated charcoal (MDAC, 0.5 g/kg every 2--4 hours), beta-blockers for hemodynamically significant tachycardia, and hemodialysis for severe toxicity. Prognosis is favorable with appropriate treatment.

\subsection{Tricyclic Antidepressant (TCA) Overdose}
\label{sec:TCA Overdose}
TCA overdose (amitriptyline, nortriptyline, imipramine, desipramine, doxepin) is a highly lethal pharmaceutical overdose due to cardiovascular and neurological toxicity. The condition results from sodium channel blockade (class IA antiarrhythmic effect) slowing Phase 0 depolarization and prolonging QRS duration, alpha-adrenergic blockade causing hypotension, anticholinergic effects, and GABA\textsubscript{A} antagonism lowering seizure threshold. Clinical features include altered mental status (often rapid onset coma), anticholinergic toxidrome, seizures, and cardiac arrhythmias, with the most specific finding being QRS widening >100 ms. Diagnosis is clinical and by ECG. Management includes airway protection, sodium bicarbonate (1--2 mEq/kg IV bolus) for QRS >100 ms or hemodynamically significant arrhythmias, benzodiazepines for seizures, and vasopressors (norepinephrine preferred) for refractory hypotension. Prognosis is favorable with early aggressive treatment.

